

\section{Parlons du Web et du REST}



\begin{frame}{Ce qu'il faut savoir (pour nous)}


Web = immense base de données documentaire! Les notions
essentielles.

\begin{redbullet}

  \item Le Web est constitué de \alert{ressources}.

    Une ressource est une entité fournissant des services et communiquant
    par message.
    
   
   \item Les ressources sont identifiées et accessibles à des \alert{URL}.
   
   \url{https://www.example.com:443/chemin/vers/doc?nom=b3d&type=json\#fragment}


   \item Les messages sont codés selon un \alert{protocole}, HTTP.
   
   
   Très utile (indispensable?): l'outil cURL pour parler le HTTP

   \item Un "message" est une enveloppe dont le contenu est un \alert{document}.
   
   Souvent le document est en HTML,  mais cela peut être un \alert{document structuré}.

\end{redbullet}

\end{frame}


\begin{frame}[fragile]
\frametitle{L'architecture REST}

REST = une définition de services Web basée sur les standards du Web.

\begin{redbullet}
  \item on s'adresse à des ressources qui fournissent des services;
  \item les messages se font en HTTP, et sont (le plus souvent) \alert{structurés} en XML ou JSON;
  \item les documents structurés fournis par une ressource sont (le plus souvent)
     produits \alert{à la volée} (par \alert{calcul}: distinction entre service et ressource statique).
\end{redbullet}

Très répandu! Excellent moyen de récupérer des informations directement exploitables
par une machine (pas comme HTML).
\end{frame}

\begin{frame}{Ressources et opérations}


\figSlideWithSize{REST}{12cm}

\vfill

Quatre opérations (celles de HTTP).
\begin{description}
  \item[\texttt{GET}] \alert{lit} (la représentation d')une ressource.
  \item[\texttt{PUT}] \alert{crée} une nouvelle ressource (ou, plus flexible: la remplace).
  \item[\texttt{POST}] envoie un message (\alert{demande de service}) à une ressource.
  \item[\texttt{DELETE}] \alert{détruit} une ressource.
\end{description}
 
\end{frame}




\begin{frame}{Un peu de rigueur}

Qu'est-ce qu'un bon (un vrai) service REST?

\begin{redbullet}
  \item Des ressources clairement organisées \alert{et stables}.
  \item Des services bien conçus (simples, clairs, concis) \alert{et stables}.
  \item Respect de la sémantique des 4 opérations. Notamment:
  
  \begin{itemize}
    \item \texttt{GET}: \alert{lecture seule}, pas d'effet de bord (état des ressources inchangé).
    \item \texttt{PUT}: \alert{création}, donc idempotent (plusieurs appels = même effet qu'un seul). 
    \item \texttt{POST}: \alert{envoi de message}: s'adresse à une ressource existante (modification possible de l'état des ressources).
  \end{itemize}
\end{redbullet}

Ce sont des principes: détails sujets à de longues discussions...

\end{frame}

\begin{frame}[fragile]{Où sont les services REST?}

Ils sont partout! Beaucoup d'applications Web (Twitter, Facebook par
exemple) ont une interface HTML \alert{et} une interface REST.

\vfill

Essayons: quel temps fait-il à Paris aujourd'hui?

\begin{minted}{bash}
curl -X GET api.openweathermap.org/data/2.5/weather?q=Paris
\end{minted}

Et à Londres? Et ailleurs?

\vfill

Vous essaierez (au moins) les services de géolocalisation de Google. 

\vfill

De quoi se constituer au fil du temps une base de documents à analyser.

\end{frame}

\section{Un système NoSQL basé sur REST: CouchDB}

\begin{frame}
\frametitle{CouchDB, en bref}

Un système typiquement NoSQL.

\begin{redbullet}
  \item gère des documents structurés (JSON);
 
  \item pas de schéma, pas de langage d'interrogation, entièrement \alert{non} normalisé.
  
  \item échanges client/serveur basés sur REST;

  \item passe en mode distribué facilement (à vérifier).
  
  \item applique quelques techniques courantes de calcul distribué (MapReduce).
\end{redbullet}

\vfill

Pas le système le plus répandu, mais facile à mettre en œuvre, intéressant à étudier.

\end{frame}


\begin{frame}{Serveurs, bases, documents}

Souvenez-vous: on s'adresse à des \alert{ressources}
\vfill

\figSlideWithSize{CouchDB}{10cm}

\end{frame}

\begin{frame}[fragile]
\frametitle{Démonstration (REST)}

Vérifions que le serveur (ici, machine locale) est prêt à nous parler.

\begin{minted}{text}
$ curl -X GET http://localhost:5984
{"couchdb":"Welcome","version":"1.0.1"}
\end{minted}

\vfill
Création d'une base de données = \texttt{PUT} d'une nouvelle ressource. 

\begin{minted}{text}
$ curl -X PUT http://localhost:5984/films
\end{minted}

\vfill 

Récupérons la nouvelle ressource (sa représentation) avec \texttt{GET}.

Création d'une base de données = \texttt{PUT} d'une nouvelle ressource. 

\begin{minted}{text}
$ curl -X GET http://localhost:5984/films
\end{minted}

\vfill

Création d'un document = \texttt{PUT} d'une nouvelle ressource dans la base de données. 

\begin{minted}{text}
$ curl -X PUT http://localhost:5984/films/doc1 
      -d '{"clef": "valeur"}'
{"ok":true,"id":"doc1","rev":"1-25eca"}
\end{minted}
\vfill

Notez le champ \texttt{rev}. Et bien sûr on peut lire la nouvelle ressource.

\begin{minted}{text}
$ curl -X GET http://localhost:5984/films/doc1
\end{minted}

\end{frame}


\begin{frame}[fragile]{Ajout de documents complets}

Envoi de messages \texttt{PUT} 
(à la nouvelle ressource) ou \texttt{POST} (à la base
de données).

\vfill

On adresse un \texttt{PUT} à une \alert{nouvelle} ressource
(il faut donc produire l'URL).

\begin{minted}{bash}
curl -X PUT http://localhost:5984/films/us 
         -d @movie_52.json -H "Content-Type: application/json" 
\end{minted}



\vfill

On adresse un \texttt{POST} à la base de données.

\begin{minted}{bash}
$ curl -X POST  http://localhost:5984/films/ -d @movie_52.json \
      -H "Content-Type: application/json"
\end{minted}

Réponse: 

\begin{minted}{json}
     {
      "ok":true,
      "id":"movie:52",
      "rev":"1-68d58b7e3904f702a75e0538d1c3015d"
     }
\end{minted}


\end{frame}


\begin{frame}[fragile]
\frametitle{Mise à jour}

CouchDB est un système \alert{multiversions}

\vfill

Mettre à jour un document = ajouter une version à une version existante;


\begin{minted}{bash}
$ curl -X PUT http://localhost:5984/films/us -d @newDoc.json 
        -H "Content-Type: image/jpg" 
{"ok":true,"id":"tsn","rev":"2-26863"}
\end{minted}

\vfill
Détruire avec \texttt{DELETE} = ajouter une version marquée "détruite".

\begin{minted}{bash}
$ curl -X DELETE http://localhost:5984/films/us?rev=2-26863
{"ok":true,"id":"tsn","rev":"3-48e92b"}
\end{minted}

\end{frame}


\begin{frame}{Bilan}

À retenir:

\begin{redbullet}
  \item REST = échange de documents sur le Web.
  \item Possibilité de constituer à peu de frais une base documentaire.
  \item  Opérations de type "dictionnaire": le minimum syndical des bases NoSQL.
  \item CouchDB: un système original concrétisant une vision de Web comme
      un serveur de documents structurés.
      
\end{redbullet}

À venir: l'interface REST de ElasticSearch, bien plus riche.

\vfill

À vous de jouer: reproduire les commandes CouchDB, trouver et interroger des services REST
en JSON ou XML.



\end{frame}


\end{document}

\begin{frame}[fragile]
\frametitle{Views in \couchdb/}

A \alert{view} is the result of a \mapreduce/ job = a list of $(key, value)$
pairs.

\vfill

Views are \alert{materialized} and indexed on the key by a B+tree.

\vfill
A \map/ function


\begin{lstlisting}[language=JavaScript]
function(doc)
{
   emit(doc.title, doc.director)
}
\end{lstlisting}

\vfill
A \reduce/ function


\begin{lstlisting}[language=JavaScript]
function (key, values) {
    return values.length;
}
\end{lstlisting}
\end{frame}

\begin{frame}[fragile]
\frametitle{Accessing views}

Here is a view (without reduce function).

\begin{lstlisting}[language=JavaScript]
function(doc)
{           
   for (i in doc.actors) {    
      actor = doc.actors[i];
      emit({"fn": actor.first_name, "ln": actor.last_name}, doc.title) ;         
   }           
}
\end{lstlisting}

\vfill

Save it in the \alert{design document} named \texttt{examples}, and name the
view \texttt{actors}. The view can be queried with:

\begin{minted}{text}
$ curl IP/movies/_design/examples/_view/actors
{"total_rows":16,"offset":0,
"rows":[
 {"id":"bed7",
   "key":{"fn":"Andrew","ln":"Garfield"},"value":"The Social network"},
 {"id":"91631b",
   "key":{"fn":"Clint","ln":"Eastwood"},"value":"Unforgiven"},
   ...
\end{minted}

\end{frame}


\begin{frame}[fragile]
\frametitle{Querying views}

A view is a B+tree index. So:

\begin{lstlisting}[language=JavaScript]
function(doc)
{           
  emit(doc.genre, doc.title) ;   
}           
\end{lstlisting}
is equivalent to 

\begin{lstlisting}[language=SQL]
create index on movies (genre);
\end{lstlisting}

\vfill

Recall the B+trees support key and range queries:


\begin{minted}{text}
$ curl $IP/movies/_design/examples/_view/genre?key=\"Drama\"
{"total_rows":5,"offset":2,"rows":[
{"id":"9163", "key":"Drama","value":"Marie Antoinette"},
{"id":"bed7", "key":"Drama","value":"The Social network"}
]}
\end{minted}

For range queries, send the two parameters \texttt{startkey}
and \texttt{endkey}.

\end{frame}

\section{Distribution}

\begin{frame}[fragile]
\frametitle{The replication primitive}

\couchdb/ supports natively one-way \alert{replication} from one instance to
another.


\begin{minted}{text}
curl -X POST $COUCHADDRESS/_replicate \
   -d '{"source": "movies",  "target": "backup", 
        "continuous": true}' \
   -H "Content-Type: application/json"
\end{minted}

\vfill

That's all: any change in \texttt{movies} is automatically reported in
\texttt{backup}.

\end{frame}

